%%%%%%%%%%%%
%
% $Autor: Wings $
% $Datum: 2019-03-05 08:03:15Z $
% $Pfad: TemplateSensor $
% $Version: 4250 $
% !TeX spellcheck = en_GB/de_DE
% !TeX encoding = utf8
% !TeX root = filename 
% !TeX TXS-program:bibliography = txs:///biber
%
%%%%%%%%%%%%

\chapter{Objektorientierte Sensorarchitektur}

Dieses Kapitel beschreibt die objektorientierte Softwarestruktur, die als gemeinsame Grundlage für die Kommunikation und Verwaltung verschiedener Sensortypen dient. Durch den Einsatz eines universellen Interfaces können unterschiedliche Sensoren – beispielsweise Strom-, Kraft-, oder Temperatursensoren – über ein einheitliches Schema angesprochen, initialisiert und ausgelesen werden. Dies ermöglicht eine modulare, wiederverwendbare und leicht erweiterbare Architektur, die unabhängig von den jeweiligen Sensortreibern funktioniert.
	
\section{Konzept der gemeinsamen Schnittstelle}
	
Das Interface \PYTHON{ISensor} definiert die grundlegenden Methoden, die jeder Sensortyp implementieren muss: 

\PYTHON{init()} zur Initialisierung, \PYTHON{update()} zur Messwertaktualisierung und \PYTHON{getValue()} zur Bereitstellung des aktuellen Messwerts. Diese Vereinheitlichung ermöglicht es, alle Sensoren im System polymorph zu behandeln – also mit derselben Schnittstelle zu arbeiten, unabhängig vom Sensortyp. 
	
\begin{itemize}
	\item \textbf{Flexibilität:} Neue Sensoren können hinzugefügt werden, ohne den bestehenden Code zu verändern.
		\item \textbf{Modularität:} Jeder Sensor ist in einer eigenen Klasse gekapselt.
		\item \textbf{Wiederverwendbarkeit:} Gleiche Schnittstelle, unterschiedliche Implementierung.
		\item \textbf{Erweiterbarkeit:} Einfache Integration neuer physikalischer Messgrö{\ss}en (z. B. Spannung, Kraft, Druck).
\end{itemize}
	
\section{Interface für die Sensoren}

	{
		\captionof{code}{Interface \PYTHON{ISensor}: gemeinsame Schnittstelle für verschiedene Sensortypen}
		\ArduinoExternal{}{../../Code/MKRWAN1310/SoftwareArchitecture/ISensor.h}
	}

	Ein Hauptprogramm kann alle Sensorobjekte in einer gemeinsamen Liste verwalten, 
	ohne deren konkreten Typ kennen zu müssen. Dadurch können beliebige Kombinationen von Sensoren dynamisch eingebunden werden:
	
	\begin{lstlisting}[language=C++,caption={Beispielhafter Polymorphismus für mehrere Sensortypen}]
		ISensor* sensors[] = {
			new CTemperatureSensor(7, "DS18B20"),
			new CCurrentSensor(A1, 10.1, 1480, "SCT013")
		};
		
		for (auto* s : sensors) {
			s->init();
			s->update();
			Serial.println(s->getValue());
		}
	\end{lstlisting}
	